\section{Background and Related Work}\label{sec:background}
\subsection{What feature-space interpretability can establish}
The feature-space interpretability survey~\citep{fsi} organizes a heterogeneous literature around method family, locus of analysis, and concept provenance. It surveys 135 studies, including 77 preprints, and builds an evidence map of 24 statements. Hanqi Jiang is first author; Lei Xing is corresponding author. The survey provides a common vocabulary for separating what can be \emph{read} from a model from what the model actually \emph{uses}.

In a probing study a researcher fits a predictor to a representation and measures whether a clinical or acquisition property is recoverable. This can establish that information is accessible to the probe under a specified data distribution. It cannot by itself establish that a downstream model relies on the same direction or that a perturbation in that direction changes a clinical answer. Other approaches seek concepts through attribution, learned dictionaries, or interventions. Their evidential force depends on controls, representation location, concept definitions, and the behavior being tested. The survey's taxonomy makes these methodological choices visible rather than treating all interpretable features as equivalent.

Four observations motivate this dissertation. First, the literature more often measures what representations encode than what models use. Second, site and scanner are the most consistently recovered properties, a reminder that strong signals may reflect acquisition rather than pathology. Third, within the survey corpus there is no controlled evidence identifying an effect of medical pre-training on interpretability; this is a statement about available evidence, not proof that pre-training has no effect. Fourth, the multimodal connector and text encoder are almost unmeasured. These omissions matter because a VLM answer depends on the path from visual representation through connector to language generation, not solely on its image encoder. Studies have found demographic information recoverable from medical images and foundation-model representations~\citep{gichoya2022ai,zheng2024demographic}, illustrating why decodability requires careful interpretation.

The distinction between clinical concepts and acquisition characteristics also complicates interpretation of saliency. A region of high importance may coincide with anatomy while a model still depends on correlated artifacts. Clinician gaze provides a possible source of region-level supervision, but gaze reflects a reading process and is not a complete causal explanation. Report sections provide another form of supervision, with findings describing observations and impressions summarizing them. The alignment work treats these clinical signals as inductive biases to be tested through downstream performance, rather than as definitive proof of reasoning.

\subsection{Related approaches to efficient medical and visual models}
The efficiency literature spans distinct resource constraints. Spiking neural networks process events with sparse operations and offer a route to lower energy, while parameter sharing can reduce duplication in multimodal generation. Dynamic querying and token selection reduce the visual information consumed during reasoning; label-efficient learning attempts to use unlabeled or normally appearing images where expert annotations are scarce. These approaches differ in measurement: token counts, architectural sharing, estimated energy, and annotation demand cannot be directly compared. EG-SpikeFormer, HybridMED, ViThinker, and PER instantiate these respective ideas~\citep{egspike,hybridmed,vithinker,per}.

The gap is a unified account of \emph{evidence use and efficiency} across a model's training, inference, and evaluation. Alignment offers relevant representations, and selective perception governs what is acquired. Audits test responses to compromised evidence, interventions examine control, and explicit readouts or candidates connect evidence to answers. None alone establishes deployment readiness; the tests identify which evidence-use properties were examined and which resource costs were measured for a given system.
